\documentclass[12pt]{article}
\usepackage[utf8]{inputenc}
\usepackage[T1]{fontenc}
\usepackage{amsmath,amssymb}
\usepackage{geometry}
\geometry{margin=2.5cm}

\begin{document}

\noindent $N=\{e,\sigma,\sigma^2\} \leq S_3$.

\bigskip
\noindent $H=\{e,\tau\}$ \quad $\tau = \begin{pmatrix}1&2&3\\2&1&3\end{pmatrix}$, \ $\tau^2=e$.

\noindent $\Rightarrow H\leq S_3$.

\bigskip
\noindent $N \trianglelefteq S_3$ \quad but \quad $H\not\trianglelefteq S_3$.

\bigskip
\noindent $a\in S_3$ and $x\in H \overset{?}{\Rightarrow} axa^{-1}\in H$.\footnote{the original reads ``$x\in N$'', but the computation that follows tests the normality of $H$ (using $\tau\in H$) and produces a result that is not in $H$ -- corrected here to $H$ for consistency with the rest of the argument.}

\[
\begin{pmatrix}1&2&3\\2&3&1\end{pmatrix}\begin{pmatrix}1&2&3\\2&1&3\end{pmatrix}\begin{pmatrix}1&2&3\\3&1&2\end{pmatrix} = \begin{pmatrix}1&2&3\\1&3&2\end{pmatrix} \notin H.
\]

\bigskip
\noindent\underline{Ex} \ 1) Let $(G,\cdot,e)$ be a group.

\noindent $\{e\}\trianglelefteq G$. \quad and \quad $G\trianglelefteq G$.

\bigskip
\noindent (2).\ $GL_2(\mathbb{R}) = \{A\in M_2(\mathbb{R}) \mid \det(A)\neq 0\}$.

\noindent $SL_2(\mathbb{R}) = \{A\in M_2(\mathbb{R}) \mid \det(A)=1\}$.

\noindent $(M_2(\mathbb{R}),\cdot)$ is a monoid.

\noindent $GL_2(\mathbb{R})$ is a stable subset of $M_2(\mathbb{R})$ with respect to $\cdot$

\noindent $(GL_2(\mathbb{R}),\cdot)$ is a non-commutative group.

\noindent $SL_2(\mathbb{R}) \trianglelefteq GL_2(\mathbb{R})$.

\bigskip
\noindent N.\ Becheanu, C.\ Vraciu -- \emph{P.S.\ de teoria grupurilor} [Group Theory Problems]. University Press.\footnote{bibliographic reference; the abbreviation ``P.S.'' is unclear in the original -- it may stand for ``Probleme [Selectate/de Seminar]'' (Selected/Seminar Problems). Title left in the original Romanian, as is standard citation practice.}

\bigskip
\noindent\underline{Lemma}: Let $(G,\cdot)$ be a group, and $R\subseteq G\times G$ an equivalence relation on $G$.

\noindent We show the following are equivalent.

\begin{enumerate}
\item[(1)] $R$ is compatible with $\cdot$. \ $\big(xRx'$ and $yRy' \Rightarrow xyRx'y'\big)$.
\item[(2)] $R$ is compatible on the left and on the right with $\cdot$

\noindent \big($xRy \Rightarrow zxRzy$ and $xzRyz$, $(\forall)\,z\in G$\big).
\end{enumerate}

\bigskip
\noindent\underline{Proof} $(1)\Rightarrow(2)$.\ $zRz \overset{(1)}{\Longrightarrow} (xzRyz)\wedge(zxRzy)$.

\bigskip
\noindent $(2) \Rightarrow (1)$. \quad $xRx' \Rightarrow xy'Rx'y'$

\noindent $yRy' \Rightarrow xyRxy'$

\noindent $\Rightarrow xyRx'y'$.

\end{document}
